% !TEX root = ../main.tex
\chapter{The World that Bagehot Knew}\label{ch:bagehot}
\begin{paracol}{2}

\begin{sources}
\begin{tabularx}{\linewidth}{@{}l l L@{}}
\toprule
\textbf{Segment} & \textbf{Length} & \textbf{Content}\\
\midrule
\seg{9}{1}{L9.1 FT: Depreciation of Iran's Currency} & 3:37 & Official and black-market rial; sanctions on the payment system.\\
\seg{9}{2}{L9.2 Reading: John Hicks} & 3:38 & Hicks, \emph{A Market Theory of Money} (1989): banks as dealers.\\
\seg{9}{3}{L9.3 Bagehot's World, Wholesale Money Market} & 7:57 & Bills of exchange and their discount.\\
\seg{9}{4}{L9.4 Economizing on Notes: Deposits, Acceptances} & 8:28 & Discounting for deposits; acceptances as guarantees.\\
\seg{9}{5}{L9.5 Managing Cash Flow: Discount, Rediscount} & 7:29 & The discount rate as a tool of cash management; rediscount between banks.\\
\seg{9}{6}{L9.6 Market Rate of Interest} & 2:44 & Decentralised management produces one market rate, the world rate.\\
\seg{9}{7}{L9.7 Central Bank and Bank Rate} & 8:13 & The Bank of England's balance sheet; bank rate.\\
\seg{9}{8}{L9.8 The Bagehot Rule, Origin of Monetary Policy} & 4:38 & Lend freely at a high rate; from lender of last resort to monetary policy.\\
\seg{9}{9}{L9.9 Limits on Central Banking: Internal vs.\ External Drain} & 10:20 & Discipline from above: gold, cooperation among central banks.\\
\bottomrule
\end{tabularx}

\smallskip
\textbf{Files.} Transcripts \file{materials/transcripts/L09-P*.txt}.
\textbf{Readings.} Bagehot, \emph{Lombard Street} (1873), two chapters; John Hicks, \emph{A Market Theory of Money} (1989), excerpt (not among the
course files).
\end{sources}

\section*{Motivation}
So far the course has dealt with \emph{par}: payments, in which a deposit or reserves are delivered to settle a debt; the interest rate stayed in the
background \ts{9}{2}{1:51}. The next two weeks ask where the interest rate comes from and what the central bank has to do with it; after the
midterm the course moves up to the exchange rate and out to asset prices, the price of risk \ts{9}{2}{2:11}. The shift is from payments to
\emph{market making}: banks, central banks and security dealers as dealers who create liquidity by standing ready to buy and sell \ts{9}{2}{2:53}.
The lecture starts in the 19th century, where things are simpler: the London money market of Bagehot \ts{9}{3}{0:30}.

% ------------------------------------------------------------------
\section{FT: depreciation of Iran's currency}\label{sec:bg-ft}

\begin{casestudy}[FT, October 2012: ``Iran uses force to strengthen rial'']
The rial has depreciated dramatically against the dollar in Tehran's bazaar, where private currency dealers trade. The government steps in and buys dollars at
28{,}000 rials; on the illegal black market a dollar costs 30{,}000, a depreciation of about 6\% \ts{9}{1}{0:23}. The
FT attributes it to Western sanctions: EU oil sanctions shrink Iran's oil revenue, and US financial sanctions make it ``slow and expensive'' for Iran to receive
payments, which are in dollars; investors stay away from the rial and citizens dump it for foreign notes \ts{9}{1}{2:07}.
\end{casestudy}
Two points for the course \ts{9}{1}{1:46}:
\begin{enumerate}
  \item the \emph{foreign exchange dealer}, a first example of the dealers studied from now on (foreign exchange itself comes after the midterm);
  \item the cause: an attempt to keep Iran out of the world \emph{payment system}, so that it cannot get dollars.
\end{enumerate}

% ------------------------------------------------------------------
\section{Reading: John Hicks}\label{sec:bg-hicks}

\begin{reading}[Hicks, \emph{A Market Theory of Money} (1989)]
Hicks's last book, written as an old man, influenced Mehrling deeply: he read it as a young professor (he started teaching at Barnard in the fall of 1987) and has
had it on the reading list since he first taught money and banking in 1996 \ts{9}{2}{0:07}. As a young man, Hicks had sent monetary economics down what
Mehrling sees as a misleading path, with an article of 1933; in 1989 he proposed another way, a market theory of money in which banks are \emph{dealers}, which is
exactly what the course does in the next two weeks. ``I think of him as the Moses of modern economics'': he saw the promised land but did not live to enter it
\ts{9}{2}{0:49}.
\end{reading}\begin{coursenote}[Which article?]
The early article is presumably Hicks's ``A Suggestion for Simplifying the Theory of Money'' (\emph{Economica}, 1935), which treated the demand for money as a choice
among assets, an approach that led to portfolio theories of money; in class the date is given as 1933 \ts{9}{2}{0:49}. Hicks received the Nobel prize in 1972.
\end{coursenote}

% ------------------------------------------------------------------
\section{Bagehot's world: bills of exchange}\label{sec:bg-bills}

Firms in a production chain (in Britain, say, a wool producer and a clothing firm) buy and sell intermediate goods from each other: business to business, the
wholesale money market \ts{9}{3}{0:51}. Firm A buys goods from firm B on credit, issuing a \emph{bill of exchange}: a promise to
pay in 90 days that refers specifically to these goods, which will have been sold by then \ts{9}{3}{2:25}. Business-to-business
trade still works this way \ts{9}{3}{3:28}.

\begin{definition}[New concepts: bill of exchange, discount]\label{def:bg-bill}
A \emph{bill of exchange} is a written order by which a debtor (here firm A) promises to pay a sum on a given date, typically 90 days, for goods received: a form of trade
credit, transferable, ``self-liquidating'' because the sale of the goods provides the means to pay. To \emph{discount} a bill is to buy it before maturity for less than
its face value; the difference is the interest (\cref{sec:st-fed}).
\end{definition}

\paragraph{The four steps} \ts{9}{3}{3:49}--\ts{9}{3}{6:23}:
\begin{enumerate}
  \item A buys goods from B with a bill;
  \item B, needing cash, discounts the bill at a bank, which pays in bank notes, say 95 for a bill of 100 (an exaggeration: it would be closer to 100);
  \item in 90 days A sells the goods to its customers for notes;
  \item A uses the notes to redeem the bill at the bank.
\end{enumerate}
\begin{taccts}
\textit{(1)}\quad\tacct[2.0cm]{Firm A}{$+$Goods}{$+$Bill}\quad
\tacct[2.0cm]{Firm B}{$-$Goods\\$+$Bill}{}

\medskip
\textit{(2)}\quad\tacct[2.0cm]{Firm B}{$-$Bill\\$+$Notes}{}\quad
\tacct[2.0cm]{Bank}{$+$Bill\\$-$Notes}{}

\medskip
\textit{(3--4)}\quad\tacct[2.0cm]{Firm A}{$-$Goods\\$+$Notes, then $-$Notes}{$-$Bill}\quad
\tacct[2.0cm]{Bank}{$-$Bill\\$+$Notes}{}
\end{taccts}
After step 2, B has been paid and is out; A owes the bank, not B \ts{9}{3}{4:31}. This is the basic discount mechanism by which firms pay each
other with the help of a bank \ts{9}{3}{6:23}.

\paragraph{The bank's opportunity and danger.} The bank gives up notes, which pay no interest, for a bill, which does: good. But it gives up a means of payment for an asset
that turns into cash only in 90 days, hopefully: a liquidity risk. Its notes are its reserves against its deposits, which can be withdrawn or transferred at any time
\ts{9}{3}{7:04}.

\section{Economizing on notes: deposits and acceptances}\label{sec:bg-economizing}

Two alternatives to step 2 save notes \ts{9}{4}{0:09}:
\begin{enumerate}[(a)]
  \item \textbf{Discount for deposits.} If B accepts a deposit instead of notes, the bank expands its balance sheet and keeps its notes. Still a liquidity risk: B discounted
        because it wants to spend, and if it pays a customer of the same bank, fine, but if it pays someone in France the bank must find notes, and behind them gold
        \ts{9}{4}{0:53}.
  \item \textbf{Acceptance.} B may prefer to hold the bill to maturity and earn the full 100 rather than 95, but wants to be able to sell it if needed. So it has the bank
        \emph{endorse} the bill: the bank promises to pay if A does not. Then the bill is ``a 90-day bill on this bank'', almost a time deposit, and easy to sell
        \ts{9}{4}{2:57}.
\end{enumerate}
\begin{taccts}
\textit{(a)}\quad\tacct[2.0cm]{Firm B}{$-$Bill\\$+$Deposit}{}\quad
\tacct[2.0cm]{Bank}{$+$Bill}{$+$Deposit of B}

\medskip
\textit{(b)}\quad\tacct[2.3cm]{Firm B}{Bill\\$+$Acceptance\\\quad(contingent)}{}\quad
\tacct[2.3cm]{Bank}{}{$+$Acceptance\\\quad(contingent)}
\end{taccts}

\begin{definition}[New concepts: acceptance, contingent liability]
A \emph{banker's acceptance} is a bank's guarantee written on a bill (physically, the word ``accepted'' and the banker's signature across it): the bank will pay if the drawee does
not. It is a \emph{contingent} liability of the bank and a contingent asset of the holder: it becomes an actual payment only if the event (the drawee's failure to pay)
occurs \ts{9}{4}{4:22}.
\end{definition}

Mehrling writes the acceptance as a separate balance sheet entry, because in the modern world it is a separate instrument: ``an early form of a credit default swap: if they
don't pay, I'll pay'' \ts{9}{4}{5:26}. The bank charges a fee, like an insurance premium, smaller than the discount, say 1\% instead of 5\%, since it is a guarantee and
not a 90-day investment; if A pays, the bank has ``money for nothing'' \ts{9}{4}{6:50}. Bankers know their firms; the typical problem is not outright default
but \emph{delay}: the goods are not sold in time, so payment is late, a liquidity issue \ts{9}{4}{7:52}.

% ------------------------------------------------------------------
\section{Managing cash flow: discount and rediscount}\label{sec:bg-cashflow}

Focus on the bank: in step 2 bills come in and notes go out, in step 4 notes come in and bills go out. ``That's the banking business'' \ts{9}{5}{0:07}. The bank holds a
\emph{ladder} of bills maturing on known dates (a 90-day bill today is an 89-day bill tomorrow), so it knows the timing of its inflows; on the other side, demand deposits that
could all go tomorrow, but whose patterns it knows from experience \ts{9}{5}{0:49}. The banker's business is matching inflows and outflows: use the notes that
come in today to discount new bills, hold as little idle cash as possible \ts{9}{5}{1:52}. Two tools:
\begin{enumerate}
  \item \textbf{The discount rate.} Short of notes, a bank gives customers ``a bad price'': it raises the rate at which it discounts, which is the same as lowering the price at which
        it buys bills. Discounting slows, maturing bills bring notes back, and later it lowers the rate again. Raising the rate shifts the balance sheet from bills into cash, lowering
        it from cash into bills \ts{9}{5}{2:33}.
  \item \textbf{Rediscount at another bank.} Instead of sending the customer down the street, the bank itself takes bills it already holds to another bank: an interbank money market
        \ts{9}{5}{5:20}.
\end{enumerate}
So each bank quotes two prices, a rate to discount for customers and a rate to rediscount with other banks: a buy--sell spread. Banks are money dealers, and the prices at which they
deal are the discount and rediscount rates \ts{9}{5}{6:46}.

\begin{intuition}[Price as a valve on flows]
Raising the discount rate does not change the bank's stock of notes at once; it changes the \emph{rate of flow}: fewer bills in, the same maturities out, so the stock of notes
recovers over the following days. The interest rate here is a control on the time derivative of reserves, set by a bank watching its own cash position.
\end{intuition}

\section{The market rate of interest}\label{sec:bg-market}

Each bank looks only at its own cash flow and moves its discount rate accordingly; the effect is \emph{coordination}. Whoever quotes the lowest rate gets all the discounts, so
rates line up, some a bit off the market in one direction or the other to influence flows: a \emph{market rate of interest} emerges \ts{9}{6}{0:07}. It moves
over time with firms' demand for discounts and the banks' decisions on deposits, acceptances and how much cash to hold: all these decisions show up in one money market rate,
uniform across the economy \ts{9}{6}{1:10}. Lombard Street was the money market not only for wool merchants but for the world: firms A and B might be in two different
countries, discounting their bills in London; the London market rate was the world rate of interest, and payment imbalances anywhere showed up in it
\ts{9}{6}{1:51}.

% ------------------------------------------------------------------
\section{The central bank and bank rate}\label{sec:bg-bankrate}

The balance sheet of the Bank of England of 9 January 1924 (after Bagehot's time, but with the structure he describes) is divided, since Peel's Act of 1844, into an Issue
Department and a Banking Department \ts{9}{7}{0:07}:
\begin{taccts}
\tacct[2.8cm]{Issue Department}{Gold coin and bullion\\Government debt (fiduciary issue)}{Notes issued}\qquad
\tacct[2.8cm]{Banking Department}{Notes (reserve)\\Government securities\\Other securities (discounts, bills)}{Deposits (public, bankers, other)}
\end{taccts}
\begin{itemize}
  \item The \textbf{Issue Department} issues the currency, mostly against gold, plus a small fiduciary issue against government debt; at the margin it can issue more notes only
        against more gold, and must pay gold to foreigners who present notes \ts{9}{7}{0:49}.
  \item The \textbf{Banking Department} holds part of the notes as its reserve (the rest circulate), government securities and other securities, discounts of trade bills; its
        liabilities are deposits, not only of banks but of firms and people: unlike the Fed, the Bank was a bank for everyone \ts{9}{7}{1:30}.
\end{itemize}
The Banking Department is in the same bill market as other banks: it discounts for the market and rediscounts for other banks \ts{9}{7}{2:31}. It quotes \emph{bank rate},
the rate at which a banker who needs notes can get them from the Bank \ts{9}{7}{3:12}.

\begin{definition}[New concepts: bank rate]
\emph{Bank rate} was the Bank of England's published minimum rate for discounting eligible bills: the central bank's rediscount rate, the price at which the banking system could
obtain the best money of the system.
\end{definition}

The Bank was privately owned and paid dividends, but it held the reserve of the whole British banking system, which was at the same time the reserve of the international money
market \ts{9}{7}{3:53}. If the banking system needs money, it discounts bills at the Bank for notes, or, in a real crisis when the Bank runs out of notes, for deposits
\ts{9}{7}{4:56}. It could behave like any bank, raising or lowering its rate to manage its own cash flow; Bagehot says it does not, because as keeper of the reserve
it has public duties, whether it admits it or not \ts{9}{7}{6:03}. Usually the market rate is below bank rate and nobody goes to the Bank; when everyone scrambles
for cash the market rate rises to bank rate, and then everyone goes to the Bank, which ``disgorges'' its notes and expands its deposits, which banks treat as cash, although deposits
are promises to pay notes and notes promises to pay gold \ts{9}{7}{7:05}.

\begin{nfigure}
\begin{tikzpicture}[font=\scriptsize,x=0.8cm,y=0.8cm]
  \draw[-{Stealth[length=3pt]}] (0,0) -- (10,0) node[right]{time};
  \draw[-{Stealth[length=3pt]}] (0,0) -- (0,4) node[above]{rate};
  \draw[very thick,thmcol] (0,3) -- (9.5,3) node[right]{bank rate};
  \draw[thick,defcol] plot[smooth] coordinates {(0,1.5) (1,1.8) (2,1.4) (3,2.0) (4,2.4) (5,3.0) (5.6,3.0) (6.2,2.6) (7,1.9) (8,1.6) (9.5,1.8)};
  \node[defcol,right] at (9.5,1.8) {market rate};
  \fill[keycol!25] (4.95,0) rectangle (5.65,3);
  \node[align=center,keycol] at (5.3,3.7) {everyone goes\\to the Bank};
\end{tikzpicture}
\caption{Bagehot's two rates \ts{9}{7}{7:05}: normally the market rate is below bank rate and nobody borrows from the Bank; in a scramble for cash the market rate rises to bank rate,
which acts as a ceiling, and the Bank lends freely there.}\label{fig:bg-rates}
\end{nfigure}

% ------------------------------------------------------------------
\section{The Bagehot rule and the origin of monetary policy}\label{sec:bg-rule}

With bank rate the Bank of England is engaged in monetary policy, whether it knows it or not: ``this is the origin of monetary policy'' \ts{9}{8}{0:08}. In its minimal
form: set a maximum rate, and when the crisis comes act passively, lending to all comers at that high rate \ts{9}{8}{0:28}.

\begin{proposition}[The Bagehot rule]\label{prop:bg-rule}
In a crisis the central bank should \emph{lend freely, at a high rate, against good security} \ts{9}{8}{0:48}. That is all Bagehot says in \emph{Lombard Street}; he does not say to move
bank rate around to influence the economy \ts{9}{8}{0:48}.
\end{proposition}

\paragraph{Elasticity and discipline in one phrase.} ``Lend freely'' is elasticity: otherwise borrowers who cannot make their payments default, and liquidity kills you quick. ``At a high
rate'' is discipline: only those who really need it come, and they repay fast, borrowing in the market again as soon as the market rate falls; the lending is self-liquidating
\ts{9}{8}{3:13}.

\paragraph{From backstop to policy.} The market knows the Bank is waiting in the wings and behaves accordingly, taking more risk, lending more; this gives the Bank influence over the
economy, and once it realises it, it can move bank rate up and down rather than leave it at, say, 6\%: lower it to encourage credit growth, raise it to discourage it (``I'm not going to help
you, you're on your own''). ``That's the art of central banking'' \ts{9}{8}{1:09}. Sitting at the top of the hierarchy, the central bank can manage it,
moving bank rate, and be lender of last resort in trouble; but it can push rates around with its balance sheet only if it lets the market move onto its balance sheet: it can relax tension or
cause it \ts{9}{8}{2:31}.

% ------------------------------------------------------------------
\section{Limits on central banking: internal and external drain}\label{sec:bg-limits}

\paragraph{Looking down: internal drain.} If notes drain out of the banks into the hands of firms, the Bank can ease the drain by releasing its notes and creating good substitutes for them,
or choose not to: a policy choice, without constraint \ts{9}{9}{0:09}.

\paragraph{Looking up: external drain.} If foreigners want notes, and ultimately gold, the Bank must give them gold. London may be the centre of the world money market, but the Bank of
England is one central bank among others, and if they will not accept deposits at the Bank as equivalent to notes and demand gold, it must pay: that is the gold standard
\ts{9}{9}{1:16}. The Bank faces the same constraint as any bank, measured in gold in and gold out for the whole country \ts{9}{9}{2:18}. And it uses the same tool: it
raises its discount rate to say ``don't come to me to discount and then take the gold to France'' \ts{9}{9}{2:38}. If it does not, it can lose all its gold and have to
suspend payments \ts{9}{9}{3:20}.

\begin{finance}[Drawing gold from the moon]
A saying of the time held that a high enough bank rate would ``draw gold from the moon''. The mechanism is a contraction of credit: at a very high rate no new bills come for discount,
maturing bills are repaid, in gold, and no gold flows out again. The Bank secures its own position and maintains the gold parity by contracting credit \ts{9}{9}{4:23}.
\end{finance}

So the Bank of England is a source of elasticity looking down, and faces discipline looking up \ts{9}{9}{5:24}. Young stresses how little gold ran the whole world monetary system: only the
gold in the Banking Department moved, the rest was locked behind the notes; this was possible because of the sophistication of the money market, where gold could be borrowed, also from other
central banks \ts{9}{9}{5:48}.

\begin{intuition}[Cooperation and its limits]
Central banks willing to borrow and lend gold among themselves could jointly relax their survival constraint, as banks do in the Fed funds market. But central bank cooperation, ``easier
said than done'', involves national interests, politicians and finance ministers. Without it the survival constraint binds on central banks, which then have no choice but to raise rates,
even if that crunches all the banks below them: discipline is transmitted down the hierarchy. ``Does this sound familiar?'' \ts{9}{9}{6:49}.
\end{intuition}

\paragraph{The money market as coordinator.} The central bank is in between, not at the absolute top: it can relax its own constraint with other central banks, and impose discipline below
only if it has discipline imposed on itself, or chooses to \ts{9}{9}{8:13}. This is the world Bagehot knew, ``not so different from our own'', and the origin of all later thinking about central
banks, before IS/LM \ts{9}{9}{8:34}. The money market rate was a reliable indicator of stress, watched by everyone, coordinating the economy and, because London was the world money
market, the world: ``enough to make you think the money market is god''. Hence \emph{Lombard Street} was a bestseller \ts{9}{9}{9:14}.

\begin{history}[Bank rate in the 19th century]
Bank rate was the Bank of England's minimum discount rate; from the 1850s, as the Bank came to manage the gold reserve actively, it was changed often, and rates of 6--10\% were used in crises
(10\% in 1857 and 1866, 9\% in 1873). Peel's Act of 1844 was suspended in 1847, 1857 and 1866, allowing the Bank to issue notes beyond the legal limit. Bagehot's \emph{Lombard Street} (1873) appeared
in the year of a further crisis.\par
\emph{References:} R.~S.~Sayers, \emph{The Bank of England 1891--1944}, Cambridge UP, 1976; Bagehot, \emph{Lombard Street}, 1873.
\end{history}

% ------------------------------------------------------------------
\flushside
\section*{Key formulas and ideas}
\begin{keyformulas}
\begin{tabularx}{\linewidth}{@{}l L@{}}
Bill of exchange & 90-day trade credit tied to goods; discounted at a bank for notes: bank gains interest, takes liquidity risk\\
Economizing on notes & discount for deposits; acceptance $=$ contingent guarantee (early CDS)\\
Discount rate & raising it slows inflow of bills, notes accumulate: a cash-flow management tool\\
Banks as money dealers & discount rate (to customers) and rediscount rate (between banks): a bid--ask spread\\
Market rate & emerges from decentralised cash management; London rate $=$ world rate\\
Bank rate & central bank's rediscount rate; normally above the market rate, a ceiling in crises\\
Bagehot rule & lend freely (elasticity) at a high rate (discipline) against good security\\
Drains & internal: policy choice; external (gold): discipline from above, relaxed only by central bank cooperation\\
\end{tabularx}
\end{keyformulas}

\section*{Exercises}

\begin{exercise}[Discount arithmetic]
A 90-day bill of 100 is discounted at a discount rate of 4\% a year (bank discount basis, 360 days).
\begin{enumerate}[(a)]
  \item What does firm B receive?
  \item What is the bank's yield on the money it lays out, expressed as simple interest on a 365-day year?
\end{enumerate}
\end{exercise}

\begin{exercise}[A bank's ladder]
A bank holds bills maturing at 10 a day for the next 90 days and notes of 50; its depositors withdraw on net 8 a day.
\begin{enumerate}[(a)]
  \item How many new bills can it discount each day without running down its notes?
  \item A large withdrawal of 40 occurs today. Describe how it can use its discount rate and the rediscount market to restore its position.
\end{enumerate}
\end{exercise}

\begin{exercise}[The two rates]
Explain, with the T-accounts of a bank and of the Bank of England, what happens when the market rate reaches bank rate and a bank rediscounts 20 of bills at the Bank, first when it is paid in
notes and then when it is paid in a deposit at the Bank.
\end{exercise}

\begin{exercise}[External drain]
Foreign central banks present notes worth 10 to the Issue Department for gold. Write the T-accounts of both departments. What can the Banking Department do with bank rate, and with what
consequences for British banks?
\end{exercise}
